\chapter{综合实验验证}
\label{chp:experiment}

本章面向完整 PlanGraph 框架开展综合实验验证，重点考察其在不同图推理基准、不同任务类型和不同输入形态下的整体性能。综合实验采用统一数据集、统一评价指标和统一运行配置，对 PlanGraph 与代表性基线进行总体比较，并进一步讨论性能收益来源、系统运行特性、失败样例分布和适用边界。

\section{实验设置}

本章实验用于在统一口径下评价完整 PlanGraph 框架的整体有效性，主要回答四个问题：PlanGraph 相比现有大语言模型图推理方法是否具有整体性能优势；在不同图推理基准测试集上的表现是否稳定；复杂图任务中的收益主要体现在哪些方面；以及性能提升是否来自规划、检索、验证和黑板调度等关键机制的协同作用。为保证不同基准和不同方法之间具有可比性，本节统一说明综合实验所采用的数据集、对比基线、运行配置和评价指标，为后续总体性能结果提供一致的实验背景。

\subsection{数据集选择}

本文选取 GraphArena、NLGraph、GraphWiz、LLM4DyG 和 GraphInstruct 五个公开图推理基准测试集进行综合评估\cite{tang2025grapharena,chen2024nlgraph,gong2024graphwiz,llm4dyg2024,graphinstruct2024}。五个数据集覆盖静态图算法、自然语言图问题、指令式图操作、组合优化和动态图推理等不同场景，能够较全面地检验 PlanGraph 在多种输入形态和任务约束下的适应性。从去除数据集内部重复任务后的任务类别看，GraphArena 包含 10 类任务，NLGraph 包含 8 类任务，GraphWiz 包含 9 类任务，LLM4DyG 包含 9 类任务，GraphInstruct 包含 21 类任务。

GraphArena 侧重抽象图算法和组合优化任务，包含邻居查询、连通分量、最短距离、图直径、图编辑距离、最大独立集、最小点覆盖、最大团、最大公共子图和旅行商问题等任务，适合考察系统在标准图结构和高组合复杂度任务中的表现。NLGraph 更强调自然语言条件下的图结构恢复和算法意图识别，覆盖连通性、最大流、环检测、拓扑排序、最短路径、二分图匹配、哈密顿路径和 GNN 推理等任务。GraphWiz 面向指令式图推理，重点考察模型能否根据明确任务指令完成算法选择和过程化求解。LLM4DyG 引入时间戳、事件序列和动态图状态变化，用于检验系统处理时间边界、状态更新和时序关系的能力。GraphInstruct 覆盖 21 类图指令跟随任务，既包含邻居查询、边存在判断、度数计算等基础操作，也包含最大流、最小生成树、欧拉路径等算法化任务。

从任务层次看，上述数据集共同覆盖三类验证目标。第一类是基础结构理解，主要检验系统能否准确恢复图结构并执行基础查询；第二类是算法化求解，主要检验规划结果能否稳定映射到图算法和程序实现；第三类是强约束或高组合复杂度任务，主要检验规划、知识检索、执行验证和反馈修复能否共同减少错误传播。由于不同基准测试集在样本规模、难度划分和答案格式方面存在差异，本文在保持公开任务定义不变的前提下统一评价指标和结果统计方式，五个基准测试集均按照公开任务划分进行评测，不对任务集合做额外抽样或裁剪。

\subsection{对比基线}

本文主要与三类代表性方法进行比较。第一类是通用强基座模型 GPT-4o 的直接推理结果，用于衡量不引入显式规划、外部知识和执行反馈时，大语言模型自身完成图推理任务的能力。第二类是代表性多智能体图推理方法 GraphTeam\cite{li2024graphteam}，用于比较角色分工式多智能体协作本身是否足以支撑复杂图任务求解。第三类是代表性程序生成式图推理方法 GCoder\cite{zhang2025gcoder}，用于比较将图问题转化为代码求解这一技术路线的效果。

选择这三类基线，是因为它们分别对应当前图推理研究中的三种典型求解方式：直接提示推理、多智能体协同和程序生成求解。GraphTeam 的公开设置以 GPT-4o 作为主要语言模型；GCoder 的公开实验报告了基于 Llama3.1-8B 和 Qwen2.5-Coder 进行 LoRA 微调后的 GCoder-L 与 GCoder-Q，本文表中 GCoder 结果取二者中的较高值。PlanGraph 与这些方法的主要区别在于，它不仅将任务转化为代码，也不仅依赖多个智能体角色分工，而是在任务理解和程序执行之间引入伪代码规划层，并通过规划引导知识检索，通过执行反馈约束程序修复和必要的回退。因此，与上述基线的比较能够说明规划反馈闭环相对于直接推理、多智能体分工和单纯程序生成路线的增益。

需要说明的是，本文没有将经典图算法或 GNN 方法作为端到端对比基线。原因在于本文任务通常以自然语言问题、图结构描述和输出格式要求共同构成输入，系统需要同时完成任务解析、约束识别、算法选择、程序生成和答案格式化。经典图算法更接近 PlanGraph 内部执行器、验证器或 oracle 的组成部分，而不是同口径的自然语言图推理系统；GNN 方法则主要面向节点分类、链路预测或图表示学习等固定监督任务，也不直接覆盖本文的完整输入输出设定。

\subsection{实现配置}

实验中，PlanGraph 使用 GPT-4o 作为基础语言模型，程序执行环境统一为 Python 3.11 与 NetworkX。为降低生成随机性带来的结果波动，主实验温度参数设置为 0.2，最大输出长度设置为 2048 tokens。系统在每个样本上只保留一条最终答案，不采用多答案投票或额外采样，以避免采样策略对总体结果造成额外影响。

外部知识检索使用固定版本的轻量图算法知识库。该知识库共包含 3349 条知识条目，覆盖最短路径、最大流、最大团、最小点覆盖、最大独立集、旅行商、哈密顿路径、连通性、拓扑排序、二分图判定、图直径、邻居查询、子图匹配和最小割等常见图任务。除第四章专门讨论的参数敏感性实验外，知识库版本、检索器和重排序规则在综合实验中保持一致。

除非特别说明，本文沿用第\ref{chp:system}章给出的默认运行配置：检索返回条数设为 3，最大局部修复轮数设为 1，最大重生成轮数设为 1，规划与检索阶段的最大尝试次数分别设为 2，单次候选程序执行超时阈值设为 15 秒，随机种子固定为 42。验证样例仅用于检查候选程序是否满足基本可执行性、任务约束和输出格式，不使用目标测试样本的标准答案；最终评测仍以公开数据集的标准答案和统一评测脚本为准。

\subsection{评价指标}

为确保不同基准测试集间结果可比较，本文在主实验中统一采用精确匹配率（Exact Match, EM）作为总体性能指标。设测试集样本数为 $N$，其中模型完全答对的样本数为 $N_c$，则精确匹配率定义为：

\begin{equation}
\mathrm{EM} = \frac{N_c}{N} \times 100\%
\end{equation}

对于代码生成类方法，精确匹配率按完全匹配口径统计。所有失败样例均保留在统计中，包括运行错误、超时错误、输出不可解析的格式错误以及程序运行成功但答案错误的情形，不进行人工过滤或事后修正。

除 EM 外，本文还在第三章和第四章使用 Hit@k、MRR、首次通过率、修复后收敛率、执行成功率、平均修复轮数和平均耗时等辅助指标。这些指标不替代总体精确匹配率，而是用于解释性能来源和系统运行特性。其中，检索质量验证使用 Hit@k 与平均倒数排名（MRR），用于衡量不同查询方式对知识命中与排序质量的影响。若第 $i$ 个查询的正确知识条目排序位置为 $\mathrm{rank}_i$，则：

\begin{equation}
\mathrm{Hit@}k = \frac{1}{N}\sum_{i=1}^{N}\mathbb{I}(\mathrm{rank}_i \le k)\times 100\%
\end{equation}

\begin{equation}
\mathrm{MRR} = \frac{1}{N}\sum_{i=1}^{N}\frac{1}{\mathrm{rank}_i}\times 100\%
\end{equation}

执行成功率用于衡量候选程序是否能够正常执行并返回可解析、可验证的结果，平均修复轮数与平均端到端耗时分别用于刻画反馈修复成本和运行代价。设可成功完成代码执行并返回可验证结果的样本数为 $N_s$，第 $i$ 个样本触发的修复轮数为 $r_i$，端到端耗时为 $t_i$，则：

\begin{equation}
\mathrm{ExecSuccess} = \frac{N_s}{N}\times 100\%
\end{equation}

\begin{equation}
\mathrm{AvgRepair} = \frac{1}{N}\sum_{i=1}^{N} r_i
\end{equation}

\begin{equation}
\mathrm{AvgTime} = \frac{1}{N}\sum_{i=1}^{N} t_i
\end{equation}

上述分项指标用于解释总体精确匹配率背后的机制来源，其中 Hit@k 与 MRR 反映检索质量，执行成功率和平均修复轮数反映反馈闭环的稳定性，平均耗时则反映完整框架的运行代价。

\section{总体性能结果}

表\ref{tab:overall_result}给出了 PlanGraph 在统一评测集合上的总体精确匹配率结果，其中任务数表示本文实际纳入主实验统计的任务类别数量。PlanGraph 在五个基准测试集中有四个取得最优结果，在其余一个基准测试集上保持竞争力，平均精确匹配率达到 96.89\%，相比最佳现有方法平均提升 2.31 个百分点。

\begin{table}[htbp]
\centering
\small
\caption{不同基准测试集上的总体精确匹配率对比（\%）}
\label{tab:overall_result}
\begin{tabular}{lcccccc}
\toprule
基准测试集 & 任务数 & GPT-4o & GraphTeam & GCoder & PlanGraph & 提升幅度 \\
\midrule
GraphArena & 10 & 49.08 & 95.14 & 94.27 & \textbf{97.74} & +2.60 \\
NLGraph & 8 & 51.40 & 97.71 & 96.90 & \textbf{98.72} & +1.01 \\
GraphWiz & 9 & 47.61 & 88.62 & 90.20 & \textbf{97.01} & +6.81 \\
LLM4DyG & 9 & 58.04 & \textbf{96.35} & 93.78 & 96.11 & -0.24 \\
GraphInstruct & 21 & 51.14 & 93.48 & 92.56 & \textbf{94.87} & +1.39 \\
\midrule
平均 & -- & 51.45 & 94.26 & 93.54 & \textbf{96.89} & +2.31 \\
\bottomrule
\end{tabular}
\end{table}

从具体结果看，PlanGraph 在 GraphWiz 上提升最为明显，较 GCoder 提升 6.81 个百分点。这意味着在算法约束较强、路径搜索更复杂的任务中，规划反馈闭环能够较明显地减少直接生成带来的错误。PlanGraph 在 GraphArena 和 NLGraph 上也取得了稳定提升，反映出本文方法对经典图算法和自然语言图问题都具备较强适配性。

在 GraphArena 上，PlanGraph 的总体精确匹配率为 97.74\%，相比 GraphTeam 和 GCoder 均有提升。该数据集同时包含基础图结构任务和组合优化任务。基础任务已经接近饱和，各方法差距较小；真正拉开差距的是图编辑距离、最大独立集、最小点覆盖、最大团、最大公共子图和旅行商问题等复杂任务。对于这类任务，模型需要同时维护候选解可行性和目标函数最优性，直接推理或一次性代码生成容易遗漏关键约束。PlanGraph 在该数据集上的提升说明，显式规划和执行验证能够在组合优化场景中发挥较稳定作用。

在 NLGraph 上，PlanGraph 取得 98.72\% 的总体精确匹配率，较最佳现有方法提升 1.01 个百分点。该数据集更强调自然语言条件下的图算法推理，任务输入往往需要先从文本中恢复图关系、任务目标和输出形式。PlanGraph 的优势在于问题智能体和规划智能体先将自然语言描述转换为结构化任务表示和算法级规划，再进入检索与编码阶段。该流程降低了自然语言表述变化对后续程序生成的影响，也使最大流、拓扑排序等依赖明确步骤的任务更容易得到稳定求解。

在 GraphWiz 上，PlanGraph 的提升最为明显，达到 97.01\%。该数据集中的哈密顿路径、最大流、最短路径和子图匹配等任务对算法步骤和输出规范要求较高。尤其是哈密顿路径类问题，若没有显式约束检查，模型常会输出看似合理但不满足路径覆盖或节点不重复要求的答案。PlanGraph 通过规划阶段显式列出路径约束，并在验证阶段检查候选程序输出是否满足这些约束，因此能够显著降低结构性错误。

在 LLM4DyG 上，PlanGraph 与最佳方法相差 0.24 个百分点，结果略低于 GraphTeam，但整体仍具有竞争力。造成这一现象的原因主要在于，LLM4DyG 不只考察静态拓扑关系，还要求模型在事件序列中处理时间边界、状态更新和关系首次成立等条件。GraphTeam 的知识库包含历史经验信息，因此在部分动态图任务上可以利用相似任务的经验解法。相比之下，PlanGraph 当前主要依赖通用 Python 图计算函数 API 与算法文档，不预先构建面向各类任务的人类专家答案或历史经验库。因而，PlanGraph 在 LLM4DyG 上略低于 GraphTeam 的同时，也说明本文方法是在较轻知识先验条件下取得了接近最优的结果，具有较低的知识构建成本和更好的跨任务适配潜力。

在 GraphInstruct 上，PlanGraph 的总体精确匹配率为 94.87\%，优于 GraphTeam 和 GCoder。该数据集的特点是任务类别较多，且许多任务表面上属于基础图操作，但输出格式和指令跟随要求较严格。PlanGraph 在邻居查询、边存在判断、连通性、度数计算、PageRank、最大流、最小生成树和欧拉路径等任务上表现稳定，说明本文方法不只适用于高复杂度组合优化任务，也能够通过结构化解析和输出格式约束提高基础图指令任务的稳健性。少数任务如 BFS、拓扑排序和公共邻居等仍存在轻微差距，主要与输出顺序、并列解表示和格式严格匹配有关。

从不同基线的表现差异来看，GPT-4o 直接推理在各个基准测试集上均明显低于程序生成式方法和多智能体方法，说明图推理任务较依赖外部结构化求解能力。GraphTeam 和 GCoder 已经取得较高性能，而 PlanGraph 除 LLM4DyG 外仍能在其余四个基准测试集上继续提升结果，表明仅依赖多智能体分工或代码生成并不充分。对于复杂图任务，规划、知识检索和执行验证之间是否形成有效配合，同样会影响最终求解质量。

进一步看，PlanGraph 的提升幅度在不同数据集之间并不完全一致。这种差异与各数据集的任务构成有关：当任务包含更多显式约束、搜索空间较大或容易出现程序实现偏差时，规划反馈闭环的收益更明显；当任务主要是基础结构查询或已有方法已经接近满分时，PlanGraph 的提升空间相对有限；当任务需要细粒度时序状态建模时，当前方法虽然保持竞争力，但仍可能受到规划粒度和验证样例覆盖不足的限制。这一结果与第\ref{chp:method}章的复杂度分层分析和第\ref{chp:system}章的协同推理机制验证能够相互印证。

\section{性能收益来源分析}

总体性能结果说明了 PlanGraph 在统一评测口径下的整体优势，但平均精确匹配率本身并不能充分解释性能提升来自哪些机制。为进一步分析收益来源，本文将实验结果与方法机制、协同推理机制和完整框架三个层面的验证目标对应起来，如表\ref{tab:experiment_mapping}所示。

\begin{table}[htbp]
\centering
\small
\caption{综合实验组织与验证目标}
\label{tab:experiment_mapping}
\begin{tabularx}{\textwidth}{p{3.2cm}p{3.0cm}X}
\toprule
实验内容 & 验证层面 & 验证目标 \\
\midrule
任务类别、NP 完全和样本难度分层分析 & 方法层 & 验证规划反馈方法在复杂、高约束图任务中的收益来源 \\
逐任务结果分析 & 方法层 & 展示方法在不同任务族中的具体表现和错误边界 \\
检索质量验证 & 协同机制层 & 验证规划引导检索与结构约束标签对知识命中效果的提升 \\
稳定性、参数敏感性与运行开销 & 协同机制层 & 验证黑板调度、反馈恢复和运行配置是否支撑稳定运行 \\
鲁棒性与消融实验 & 框架层 & 验证关键模块必要性以及输入扰动下的运行稳定性 \\
总体性能对比 & 综合层 & 在统一实验口径下综合评价 PlanGraph 的整体有效性 \\
\bottomrule
\end{tabularx}
\end{table}

从表\ref{tab:experiment_mapping}可以看出，PlanGraph 的性能收益并非来自单一组件，而是来自方法层、协同机制层和框架层之间的协同。方法层实验主要说明规划反馈是否能够改善复杂图任务的求解质量；协同机制层实验主要说明多智能体协同、黑板调度、共享缓存管理和反馈恢复是否能够支撑方法稳定运行；综合层实验则在统一评价口径下检验完整框架相对于代表性基线的整体表现。

从方法层看，规划反馈的收益并不是均匀分布在所有任务上。任务类别分析说明 PlanGraph 并非只在某一类任务上偶然有效；NP 完全任务和困难样本分析说明其收益主要集中在高约束、高组合复杂度场景；逐任务结果则展示了不同数据集内部的细粒度差异。由此可以看出，规划反馈机制主要在需要显式约束维护、算法步骤拆解和执行结果校验的任务中发挥作用。

从协同机制层看，方法收益能否稳定转化为最终性能，还取决于黑板调度和反馈恢复是否能够稳定运行。重复运行稳定性对应黑板事件和共享缓存链能否抑制多轮协作中的状态漂移；检索质量对应规划引导查询是否真正改善知识命中；参数敏感性和运行开销对应调度预算是否合理；消融实验则对应规划、检索和验证模块是否具有不可替代性。这些实验共同说明，协同推理机制不是方法之外的附属实现，而是保证 PlanGraph 可追踪、可恢复和可复现的必要条件。

从完整框架看，PlanGraph 的总体表现可以理解为“规划反馈带来的求解质量提升”和“黑板调度带来的运行稳定性”共同作用的结果。前者解释复杂图任务上为什么能够减少约束遗漏和算法选择错误，后者解释多阶段求解过程为什么能够被记录、修复和复核。因此，本章的总体性能结果不是孤立的数值汇总，而是对前述方法机制和协同推理机制的统一检验。

\section{典型案例分析}

为了更直观地说明 PlanGraph 的工作过程，本文选取 GraphArena 中一个旅行商问题样例进行分析。该样例给定一个五节点带权无向图，要求从节点 A 出发，访问所有节点恰好一次并返回起点，同时使路径总代价最小。

\begin{figure}[htbp]
\centering
\includegraphics[width=0.95\textwidth]{figures/paper/case2.pdf}
\caption{案例求解流程图}
\label{fig:tsp_case_pipeline}
\end{figure}

常见错误主要包括：未访问全部节点，生成的路径并非哈密顿回路；节点出现重复访问，违反每个节点只访问一次的约束；虽然路径形式正确，但代价计算错误；以及混淆无向边和有向边，导致路径不可达。这些错误的根本原因在于，模型很难在一次生成中同时兼顾图构建、约束检查、路径枚举和最优性判定。

PlanGraph 对该任务的求解意图可以概括为：构建带权无向图，固定起点 A，枚举满足哈密顿回路约束的候选路径，计算每条路径总代价，选择总代价最小的回路并格式化输出。在此基础上，系统会围绕路径搜索、组合遍历和图构建等关键点补充相应知识，并通过测试样例检查是否恰好访问全部节点、是否回到起点、是否按权重正确计算代价以及输出格式是否符合要求。若某次生成的候选程序遗漏返回起点这一约束，验证阶段可以通过测试样例及时发现该问题，并将错误反馈给修复阶段。经过修正后，系统输出的最优路径为
\begin{equation}
A \rightarrow B \rightarrow E \rightarrow D \rightarrow C \rightarrow A
\end{equation}
对应总代价为 18。

\section{失败样例分析}

虽然 PlanGraph 在多个基准测试集上取得了较高精确匹配率，但仍然存在失败样例。对这些失败样例进行分析，可以进一步说明完整框架的适用边界，也能解释为什么某些数据集或任务类别上的提升幅度有限。结合方法层误差归因和协同过程运行日志，失败样例的主要来源可归纳为六类，如表\ref{tab:failure_summary_ch5}所示。

\begin{table}[htbp]
\centering
\small
\caption{失败样例来源与改进方向总结}
\label{tab:failure_summary_ch5}
\begin{tabularx}{\textwidth}{p{2.6cm}p{4.1cm}p{3.4cm}X}
\toprule
失败来源 & 典型表现 & 影响任务 & 改进方向 \\
\midrule
任务解析偏差 & 从自然语言中抽取的图结构、起终点或任务类型不完整 & 场景化路径任务、动态图查询 & 增强结构化解析约束，引入字段完整性检查 \\
规划粒度不足 & 规划只给出泛化步骤，未展开时间过滤、闭环约束或边界条件 & 动态图、旅行商、哈密顿路径 & 细化任务族规划模板，增加约束标签 \\
知识检索不匹配 & 检索结果语义相关但 API 输入输出或适用图类型不一致 & 最大流、匹配、子图任务 & 加强任务标签过滤和失败反馈重排序 \\
候选程序实现错误 & 参数名、边权字段、容量字段或返回值处理错误 & 最短路径、最大流、组合优化任务 & 利用执行反馈进行局部修复和补充检索 \\
验证覆盖不足 & 小样例通过，但目标样例存在特殊结构或边界情况 & 大规模图、长尾任务 & 自动生成更多边界验证样例 \\
输出格式不一致 & 答案内容正确但列表顺序、大小写或文本格式不匹配 & GraphInstruct、NLGraph & 增加格式约束和评测脚本对齐检查 \\
\bottomrule
\end{tabularx}
\end{table}

第一类失败主要来自任务解析偏差。在自然语言输入较长、边关系描述较分散或任务目标隐含在背景语句中时，问题智能体可能未能完整抽取节点集合、边集合、起点终点或关键约束。此类错误一旦发生，后续规划和程序生成都会建立在不完整的任务表示上，即使代码能够正常执行，最终答案仍可能偏离题意。对于这类问题，单纯增加修复轮数并不一定有效，更合理的方向是在解析阶段引入字段完整性检查，并在缺失关键字段时触发重新解析或要求规划阶段显式确认输入条件。

第二类失败来自规划粒度不足。这类错误在动态图和组合优化任务中较为典型。例如，动态图任务需要先根据时间窗口过滤事件，再构建有效图并执行查询；旅行商和哈密顿路径任务则需要显式检查节点覆盖、路径连通和回路闭合。如果规划只写成“求解最优路径”或“判断动态图关系”这类泛化描述，编码阶段就容易遗漏真正决定答案正确性的边界条件。该现象说明规划不应只是自然语言解释，而应具有足够清晰的算法动作和约束标签。

第三类失败与知识检索不匹配有关。图算法库中有些 API 的功能描述相近，但适用条件、参数名称和返回值形式并不完全相同。例如，最短路径、路径长度、所有简单路径和加权路径搜索都与路径任务相关，但它们的输入参数和输出结构不同；最大流任务中也存在流值、流字典、容量字段和残量网络等多个相关概念。若检索结果只在语义上相关，却没有与当前任务的图类型和约束形式对齐，编码智能体就可能选择错误接口。第\ref{chp:system}章的检索质量实验已经表明，加入规划伪代码和结构约束标签能够显著改善这一问题，但在长尾任务上仍有继续优化空间。

第四类失败来自候选程序实现错误。这类错误包括变量名不一致、边权字段读取错误、容量字段缺失、返回值解析错误和格式化输出异常等。与规划错误不同，实现错误通常可以通过执行反馈定位，因此 PlanGraph 的局部修复机制对这类失败具有较好恢复能力。表\ref{tab:stability_observation}中的修复后收敛率也说明，多数实现层错误能够在有限轮数内被纠正。剩余问题主要出现在错误同时涉及多个位置，或候选程序整体结构与规划不一致的情况，此时局部修复可能不如有界重生成有效。

第五类失败与验证覆盖不足有关。执行验证可以发现大量约束违反和运行错误，但验证样例不可能覆盖所有边界情况。若验证样例规模过小或结构过于简单，候选程序可能在小样例上通过，却在目标图上因特殊结构失败。例如，路径任务中目标图可能存在断边、重复边或不可达节点；组合优化任务中目标图规模增大后可能触发搜索超时；动态图任务中某些事件边界只有在特定时间点才会暴露错误。因此，验证样例生成质量直接影响反馈修复的效果。

第六类失败来自输出格式不一致。精确匹配评测通常对答案格式较敏感，即使系统已经计算出正确结果，如果最终输出包含多余自然语言、列表顺序不符合要求、大小写不一致或集合表示方式不同，也可能被判为失败。GraphInstruct 和 NLGraph 中部分基础图任务的差距主要与这一问题相关。该类错误并不表示系统缺乏图推理能力，但会影响最终评测结果。因此，后续可进一步加强输出格式模板和评测脚本对齐，使答案生成阶段更加稳定。

总体来看，失败样例并非集中在单一模块，而是分布在解析、规划、检索、编码、验证和输出多个阶段。这也从反面说明，图推理任务需要一个可追踪的多阶段求解框架。如果系统只输出最终答案，就很难判断失败原因；而 PlanGraph 保存的任务表示、规划、检索结果、候选代码、验证反馈和修复记录，使研究者可以将错误定位到相对明确的阶段。上述分析也提示后续改进方向应集中在三个方面：提高规划表示的约束粒度，提高检索结果与任务条件的匹配程度，以及提高验证样例对边界条件的覆盖能力。

\section{综合结果讨论}

综合前述实验结果，可以得到三个主要结论。第一，PlanGraph 的性能提升主要来自复杂图任务场景，而不是简单结构查询任务上的边际收益。在非 NP 完全任务和简单样本上，不同强方法之间差距较小；当任务进入 NP 完全或困难样本区间后，规划、检索与验证模块的协同作用更加明显。第二，PlanGraph 的优势并不只来自代码生成能力，而是来自任务解析、规划生成、知识检索、程序执行和反馈修复之间的闭环协作。消融实验显示，移除规划模块、检索模块或验证模块都会造成性能下降，其中规划和验证对结果影响尤为明显。第三，协同推理机制对方法稳定落地具有必要性。黑板事件、缓存版本、调度状态和配置管理使得多智能体协作过程可以被追踪、复现和恢复，避免系统退化为无序的多轮自由对话。

从方法层面看，PlanGraph 的核心作用可以概括为将图推理任务从一次性生成问题转化为分阶段约束求解问题。任务解析阶段降低自然语言输入的不确定性，规划阶段明确算法路线和关键约束，检索阶段补充外部图算法知识，编码阶段将规划转化为可执行程序，验证阶段检查程序是否满足任务条件，反馈阶段则根据错误类型选择局部修复、补充检索或重新规划。由于每一阶段都有明确输入和输出，系统能够避免在单一上下文中同时处理所有目标，从而降低复杂任务中的错误耦合。

从运行机制看，PlanGraph 的关键并不是简单设置多个智能体，而是通过黑板调度将多个智能体组织为可控流程。多智能体自由对话虽然灵活，但容易出现角色职责重叠、上下文覆盖和失败后无序重试等问题。本文的黑板调度通过状态转移、事件记录和缓存版本管理，使各智能体只处理本阶段任务，并把阶段结果显式写入共享状态。这样不仅提高了运行稳定性，也使失败样例能够被回放和复核。对于需要论文实验支撑的方法而言，可复现和可追踪本身就是协同推理机制的重要价值。

从实验层面看，PlanGraph 的收益具有一定条件性。它最适合结构约束清晰、算法实现复杂且结果可验证的图推理任务，例如最短路径、最大流、哈密顿路径、旅行商问题和部分动态图查询。对于简单邻居查询或边存在判断，完整规划反馈流程可能带来额外开销；对于需要复杂语义理解或开放式解释的任务，当前基于程序执行的验证方式也未必能够覆盖全部正确性要求。因此，本文并不把 PlanGraph 描述为所有图智能问题的通用求解器，而是将其定位为面向自然语言图推理与程序化图求解之间的协同框架。

从运行代价看，PlanGraph 相比直接生成需要额外的规划、检索和验证调用，这会增加端到端耗时。然而，对于复杂图任务而言，最低时延并不是唯一目标。若系统快速输出错误答案，后续仍需要人工检查或重新求解；而通过规划和验证提前发现错误，虽然增加了部分计算成本，但能显著提高结果可信度。第\ref{chp:system}章的运行开销分析也表明，最短路径等简单任务开销较低，而旅行商等组合优化任务开销较高，这与任务复杂度和验证需求是一致的。

同时，本文方法也存在一定局限。对于极简单的图查询，完整规划反馈流程会带来额外开销，未必比直接程序生成更经济；对于细粒度动态图推理或复杂时序约束任务，当前规划模板和验证样例仍可能不够充分；对于知识库未覆盖的长尾算法，检索模块提供的实现依据也可能不足。因此，后续工作可以从三个方向继续改进：其一，构建更细粒度的时序规划模板，提高动态图任务中的时间边界和状态更新处理能力；其二，利用失败反馈动态调整检索排序，提高知识条目与当前错误类型之间的匹配度；其三，扩展验证样例生成策略，使验证阶段能够覆盖更多边界条件和复杂约束。

此外，本文当前实验主要基于公开图推理基准测试集。虽然这些数据集已经覆盖抽象图问题和部分场景化任务，但真实应用中的图推理问题可能更开放，例如物流路径规划、网络运维、知识图谱问答、依赖关系分析和流程优化等。这些任务通常包含更多领域背景、隐含约束和多轮交互需求。PlanGraph 的多阶段框架为处理此类问题提供了基础，但仍需要结合领域知识库、任务适配器和更强的验证器进一步扩展。换言之，本文验证的是规划反馈式多智能体图推理框架的有效性和可行性，后续还需要在更复杂真实场景中继续检验其应用价值。

\section{本章小结}

本章对 PlanGraph 的综合实验结果进行了汇总与分析。在统一实验设置下，PlanGraph 在五个基准测试集中有四个取得最优结果，平均精确匹配率达到 96.89\%，较最佳现有方法平均提升 2.31 个百分点。本章进一步说明了五个数据集的任务覆盖范围、基线选择依据和评价口径，并结合总体性能结果分析了 PlanGraph 在不同数据集上的表现差异。

本章还从方法层、协同机制层和综合层三个角度分析了性能收益来源：方法验证部分说明规划反馈在复杂图任务中的作用，协同推理机制验证部分说明黑板调度和反馈恢复对稳定运行的支撑，综合性能结果则检验完整 PlanGraph 框架在统一评测口径下的表现。失败样例分析进一步表明，错误来源分布在任务解析、规划、检索、编码、验证和输出多个阶段；保留这些阶段结果和恢复记录，有助于定位失败原因并解释不同任务上的性能差异。
